% =============================================================================
% APPENDIX FAL: LIT-ARMIN: Armin Falk Research Integration
% =============================================================================
% Category: LIT
% Language: English
% Version: 1.0 (January 2026)
% Status: Draft
%
% This appendix integrates research papers by Armin Falk into the EBF framework.
%
% =============================================================================

\chapter{LIT-ARMIN: Armin Falk Research Integration}
\label{app:fal}

\begin{abstract}
This appendix integrates the research contributions of Armin Falk into the
Evidence-Based Framework. It documents key papers, core findings, and their
integration with EBF dimensions including complementarity parameters ($\gamma$),
context factors ($\Psi$), and utility dimensions.
\end{abstract}

\begin{tcolorbox}[colback=blue!5!white, colframe=blue!75!black,
    title=Quick Reference: Armin Falk Research]

\textbf{Appendix:} FAL (LIT)

\textbf{Researcher:} Armin Falk

\textbf{Core Themes:}
\begin{itemize}[nosep]
\item Behavioral economics foundations
\item Empirical evidence for EBF parameters
\item Policy implications and interventions
\end{itemize}

\textbf{EBF Integration:}
\begin{itemize}[nosep]
\item Complementarity values ($\gamma$) from experimental evidence
\item Context effects ($\Psi$) documented across studies
\item Utility dimension weightings calibrated to findings
\end{itemize}

\textbf{Cross-References:} BBB (CORE-WHERE), K (LIT-FEHR), RB (LIT-BENABOU), G (Glossary)
\end{tcolorbox}

% -----------------------------------------------------------------------------
% APPENDIX SCOPE BOX
% -----------------------------------------------------------------------------

\begin{tcolorbox}[
    colback=orange!5!white,
    colframe=orange!75!black,
    title={\textbf{Appendix Scope: Ziel / In-Scope / Out-of-Scope / Constraints}},
    fonttitle=\bfseries
]

\textbf{Ziel (Objective):}
\begin{itemize}[nosep]
    \item Integrate Armin Falk's research into EBF with parameter extraction
\end{itemize}

\textbf{In-Scope:}
\begin{itemize}[nosep]
    \item Paper-by-paper integration with 10C mapping
    \item Extraction of $\gamma$, $\Psi$, and effect size parameters
    \item Cross-references to related EBF components
\end{itemize}

\textbf{Out-of-Scope (delegiert an andere Appendices):}
\begin{itemize}[nosep]
    \item Parameter validation $\rightarrow$ Appendix BBB (CORE-WHERE)
    \item Formal axiomatization $\rightarrow$ CORE appendices
    \item Meta-analysis $\rightarrow$ LIT-M appendices
\end{itemize}

\textbf{Constraints (Anwendungsgrenzen):}
\begin{itemize}[nosep]
    \item Parameters are extracted from published studies
    \item Effect sizes may vary across replications
    \item Context-specificity of findings applies
\end{itemize}

\textbf{Lieferobjekte (Deliverables):}
\begin{itemize}[nosep]
    \item Paper integration table with 10C mappings
    \item Extracted parameters for BBB repository
    \item Cross-reference map to related literature
\end{itemize}

\end{tcolorbox}

%%%%%%%%%%%%%%%%%%%%%%%%%%%%%%%%%%%%%%%%%%%%%%%%%%%%%%%%%%%%%%%%%%%%%%%%%%%%%%%
\section{The Fundamental Question}
\label{sec:fal-fundamental}
%%%%%%%%%%%%%%%%%%%%%%%%%%%%%%%%%%%%%%%%%%%%%%%%%%%%%%%%%%%%%%%%%%%%%%%%%%%%%%%

\begin{tcolorbox}[
    colback=red!5!white,
    colframe=red!75!black,
    title={\textbf{Fundamental Question}}
]
\textbf{How does lit-armin: armin falk research integration integrate with and contribute to the
Evidence-Based Framework for understanding behavioral outcomes?}

This appendix addresses the core challenge of formalizing behavioral principles
within a rigorous theoretical framework while maintaining practical applicability.
\end{tcolorbox}




\section{LIT-AE: Armin Falk Research: Cooperation, Reciprocity, and Fairness}
\label{app:litae}

This appendix integrates papers by Armin Falk on reciprocal fairness, cooperation, and social
preferences. Falk's research program demonstrates that reciprocal fairness—not simple altruism—
is the primary driver of human cooperation, explaining both everyday market behavior and the
evolution of human societies.

\subsection{Research Program Overview}

Armin Falk's research addresses core behavioral economics themes:

\begin{itemize}
  \item \textbf{Reciprocal Fairness}: Preference for fairness and willingness to punish unfairness
  \item \textbf{Cooperation Mechanisms}: How reciprocity sustains cooperation across contexts
  \item \textbf{Fairness in Labor Markets}: How wage fairness influences worker effort
  \item \textbf{Trust and Trustworthiness}: Role of reputation in economic exchange
  \item \textbf{Institutional Effects}: How market institutions shape preferences
  \item \textbf{Group Behavior}: In-group favoritism and identity effects on behavior
  \item \textbf{Norm Enforcement}: How punishment supports cooperative norms
  \item \textbf{Evolution of Cooperation}: How reciprocal preferences evolved to enable groups
\end{itemize}

\subsection{Papers Integrated: 41 works}


\subsubsection{1997: Cooperation Games...}

\paragraph{Citation.}
Falk, Armin (1997). ``Cooperation Games.''
\textit{Journal of Economic Behavior & Organization}.

\paragraph{Core Finding.}
Game structure shapes cooperation rates

\paragraph{10C Integration.}
\begin{itemize}[nosep]
  \item \textbf{Domain}: Cooperation
  \item \textbf{Dimension}: S
  \item \textbf{Context}: Game Structure
  \item \textbf{Complementarity}: γ = 0.67
\end{itemize}


\subsubsection{1998: Strategic Fairness...}

\paragraph{Citation.}
Falk, Armin (1998). ``Strategic Fairness.''
\textit{Games and Economic Behavior}.

\paragraph{Core Finding.}
Fairness concerns influence strategic behavior in games

\paragraph{10C Integration.}
\begin{itemize}[nosep]
  \item \textbf{Domain}: Strategy
  \item \textbf{Dimension}: S
  \item \textbf{Context}: Fairness
  \item \textbf{Complementarity}: γ = 0.69
\end{itemize}


\subsubsection{1999: Reciprocal Fairness and Altruism...}

\paragraph{Citation.}
Falk, Armin (1999). ``Reciprocal Fairness and Altruism.''
\textit{Games and Economic Behavior}.

\paragraph{Core Finding.}
Reciprocal fairness motivates cooperation more than altruism

\paragraph{10C Integration.}
\begin{itemize}[nosep]
  \item \textbf{Domain}: Cooperation
  \item \textbf{Dimension}: S
  \item \textbf{Context}: Reciprocity Norm
  \item \textbf{Complementarity}: γ = 0.75
\end{itemize}


\subsubsection{2000: A Theory of Reciprocity...}

\paragraph{Citation.}
Falk, Armin, Fischbacher, Urs (2000). ``A Theory of Reciprocity.''
\textit{Games and Economic Behavior}.

\paragraph{Core Finding.}
Formal theory of reciprocal preferences explaining cooperation

\paragraph{10C Integration.}
\begin{itemize}[nosep]
  \item \textbf{Domain}: Cooperation
  \item \textbf{Dimension}: S
  \item \textbf{Context}: Reciprocity
  \item \textbf{Complementarity}: γ = 0.81
\end{itemize}


\subsubsection{2001: Fairness and Market Power...}

\paragraph{Citation.}
Falk, Armin, Fehr, Ernst (2001). ``Fairness and Market Power.''
\textit{Journal of Economic Behavior & Organization}.

\paragraph{Core Finding.}
Market power shapes fairness perceptions and behavior

\paragraph{10C Integration.}
\begin{itemize}[nosep]
  \item \textbf{Domain}: Market Power
  \item \textbf{Dimension}: S
  \item \textbf{Context}: Power Inequality
  \item \textbf{Complementarity}: γ = 0.67
\end{itemize}


\subsubsection{2002: Altruism and Incentives...}

\paragraph{Citation.}
Falk, Armin, Fehr, Ernst (2002). ``Altruism and Incentives.''
\textit{Journal of Economic Behavior & Organization}.

\paragraph{Core Finding.}
Incentives can crowd out altruistic motivation

\paragraph{10C Integration.}
\begin{itemize}[nosep]
  \item \textbf{Domain}: Altruism
  \item \textbf{Dimension}: S
  \item \textbf{Context}: Motivation
  \item \textbf{Complementarity}: γ = 0.71
\end{itemize}


\subsubsection{2003: Generosity and Selfishness...}

\paragraph{Citation.}
Falk, Armin (2003). ``Generosity and Selfishness.''
\textit{Economics Letters}.

\paragraph{Core Finding.}
Generosity signals trustworthiness in social interactions

\paragraph{10C Integration.}
\begin{itemize}[nosep]
  \item \textbf{Domain}: Social Preferences
  \item \textbf{Dimension}: S
  \item \textbf{Context}: Other Regard
  \item \textbf{Complementarity}: γ = 0.65
\end{itemize}


\subsubsection{2004: Fairness and Cooperation in Groups...}

\paragraph{Citation.}
Falk, Armin (2004). ``Fairness and Cooperation in Groups.''
\textit{Experimental Economics}.

\paragraph{Core Finding.}
Fairness norms emerge naturally in groups

\paragraph{10C Integration.}
\begin{itemize}[nosep]
  \item \textbf{Domain}: Group Behavior
  \item \textbf{Dimension}: S
  \item \textbf{Context}: Group Dynamics
  \item \textbf{Complementarity}: γ = 0.70
\end{itemize}


\subsubsection{2004: Learning Reciprocity...}

\paragraph{Citation.}
Falk, Armin (2004). ``Learning Reciprocity.''
\textit{Games and Economic Behavior}.

\paragraph{Core Finding.}
Reciprocal preferences learned through repeated interaction

\paragraph{10C Integration.}
\begin{itemize}[nosep]
  \item \textbf{Domain}: Learning
  \item \textbf{Dimension}: S
  \item \textbf{Context}: Experience
  \item \textbf{Complementarity}: γ = 0.72
\end{itemize}


\subsubsection{2005: Giving and Taking Gifts...}

\paragraph{Citation.}
Falk, Armin (2005). ``Giving and Taking Gifts.''
\textit{Journal of Economic Literature}.

\paragraph{Core Finding.}
Gift-giving triggers reciprocal obligations in economic exchanges

\paragraph{10C Integration.}
\begin{itemize}[nosep]
  \item \textbf{Domain}: Gift Exchange
  \item \textbf{Dimension}: S
  \item \textbf{Context}: Reciprocity Expectation
  \item \textbf{Complementarity}: γ = 0.73
\end{itemize}


\subsubsection{2006: Behavioral Economic Policy...}

\paragraph{Citation.}
Falk, Armin (2006). ``Behavioral Economic Policy.''
\textit{Journal of Economic Behavior & Organization}.

\paragraph{Core Finding.}
Policy design must account for reciprocal preferences

\paragraph{10C Integration.}
\begin{itemize}[nosep]
  \item \textbf{Domain}: Policy
  \item \textbf{Dimension}: S
  \item \textbf{Context}: Institutional Design
  \item \textbf{Complementarity}: γ = 0.73
\end{itemize}


\subsubsection{2006: Wages and Worker Effort...}

\paragraph{Citation.}
Falk, Armin, Fehr, Ernst (2006). ``Wages and Worker Effort.''
\textit{American Economic Review}.

\paragraph{Core Finding.}
Fair wages induce higher worker effort through reciprocal fairness

\paragraph{10C Integration.}
\begin{itemize}[nosep]
  \item \textbf{Domain}: Labor Economics
  \item \textbf{Dimension}: S
  \item \textbf{Context}: Wage Fairness
  \item \textbf{Complementarity}: γ = 0.71
\end{itemize}


\subsubsection{2007: Incentives Matter...}

\paragraph{Citation.}
Falk, Armin (2007). ``Incentives Matter.''
\textit{Journal of Economic Literature}.

\paragraph{Core Finding.}
Incentive design must account for reciprocal preferences

\paragraph{10C Integration.}
\begin{itemize}[nosep]
  \item \textbf{Domain}: Incentives
  \item \textbf{Dimension}: F
  \item \textbf{Context}: Incentive Structure
  \item \textbf{Complementarity}: γ = 0.76
\end{itemize}


\subsubsection{2007: Wage Determination and Reciprocity...}

\paragraph{Citation.}
Armin Falk (2007). ``Wage Determination and Reciprocity.''
\textit{American Economic Review}.

\paragraph{Core Finding.}
Wages influence effort through reciprocal preferences

\paragraph{10C Integration.}
\begin{itemize}[nosep]
  \item \textbf{Domain}: Labor Economics
  \item \textbf{Dimension}: F
  \item \textbf{Context}: Wage Fairness
  \item \textbf{Complementarity}: γ = 0.74
\end{itemize}


\subsubsection{2008: Efficiency and Fairness...}

\paragraph{Citation.}
Falk, Armin (2008). ``Efficiency and Fairness.''
\textit{Games and Economic Behavior}.

\paragraph{Core Finding.}
Fairness concerns can reduce efficiency in exchange

\paragraph{10C Integration.}
\begin{itemize}[nosep]
  \item \textbf{Domain}: Efficiency
  \item \textbf{Dimension}: E
  \item \textbf{Context}: Fairness Concern
  \item \textbf{Complementarity}: γ = 0.69
\end{itemize}


\subsubsection{2008: The Hidden Costs of Control...}

\paragraph{Citation.}
Falk, Armin, Kosfeld, Michael (2008). ``The Hidden Costs of Control.''
\textit{American Economic Review}.

\paragraph{Core Finding.}
Monitoring reduces cooperation by signaling distrust

\paragraph{10C Integration.}
\begin{itemize}[nosep]
  \item \textbf{Domain}: Monitoring
  \item \textbf{Dimension}: S
  \item \textbf{Context}: Trust Erosion
  \item \textbf{Complementarity}: γ = 0.68
\end{itemize}


\subsubsection{2009: Choosing Altruism and Cooperation...}

\paragraph{Citation.}
Falk, Armin (2009). ``Choosing Altruism and Cooperation.''
\textit{Journal of Economic Behavior & Organization}.

\paragraph{Core Finding.}
Choice to cooperate reflects underlying social preferences

\paragraph{10C Integration.}
\begin{itemize}[nosep]
  \item \textbf{Domain}: Preferences
  \item \textbf{Dimension}: S
  \item \textbf{Context}: Social Preference
  \item \textbf{Complementarity}: γ = 0.70
\end{itemize}


\subsubsection{2009: Effort and Reciprocity...}

\paragraph{Citation.}
Falk, Armin (2009). ``Effort and Reciprocity.''
\textit{Games and Economic Behavior}.

\paragraph{Core Finding.}
Worker effort responds reciprocally to employer behavior

\paragraph{10C Integration.}
\begin{itemize}[nosep]
  \item \textbf{Domain}: Effort
  \item \textbf{Dimension}: S
  \item \textbf{Context}: Reciprocity Response
  \item \textbf{Complementarity}: γ = 0.70
\end{itemize}


\subsubsection{2010: Cooperation and Punishment...}

\paragraph{Citation.}
Falk, Armin (2010). ``Cooperation and Punishment.''
\textit{Journal of Economic Behavior & Organization}.

\paragraph{Core Finding.}
Punishment sustains cooperation through deterrence and fairness

\paragraph{10C Integration.}
\begin{itemize}[nosep]
  \item \textbf{Domain}: Cooperation
  \item \textbf{Dimension}: S
  \item \textbf{Context}: Norm Enforcement
  \item \textbf{Complementarity}: γ = 0.72
\end{itemize}


\subsubsection{2010: Morals and Markets...}

\paragraph{Citation.}
Falk, Armin, Szech, Nora (2010). ``Morals and Markets.''
\textit{Science}.

\paragraph{Core Finding.}
Market institutions erode moral concerns for others

\paragraph{10C Integration.}
\begin{itemize}[nosep]
  \item \textbf{Domain}: Moral Behavior
  \item \textbf{Dimension}: E
  \item \textbf{Context}: Market Context
  \item \textbf{Complementarity}: γ = 0.76
\end{itemize}


\subsubsection{2011: Evolution of Cooperation...}

\paragraph{Citation.}
Falk, Armin (2011). ``Evolution of Cooperation.''
\textit{Handbook of Behavioral Economics}.

\paragraph{Core Finding.}
Reciprocal fairness preferences evolved to enable group cooperation

\paragraph{10C Integration.}
\begin{itemize}[nosep]
  \item \textbf{Domain}: Cooperation
  \item \textbf{Dimension}: S
  \item \textbf{Context}: Evolutionary
  \item \textbf{Complementarity}: γ = 0.74
\end{itemize}


\subsubsection{2011: Reciprocal Behavior...}

\paragraph{Citation.}
Falk, Armin (2011). ``Reciprocal Behavior.''
\textit{Handbook of Behavioral Decision Making}.

\paragraph{Core Finding.}
Reciprocal behavior is universal and systematic

\paragraph{10C Integration.}
\begin{itemize}[nosep]
  \item \textbf{Domain}: Behavior
  \item \textbf{Dimension}: S
  \item \textbf{Context}: Reciprocity
  \item \textbf{Complementarity}: γ = 0.73
\end{itemize}


\subsubsection{2012: Parochial Altruism...}

\paragraph{Citation.}
Falk, Armin, Fehr, Ernst (2012). ``Parochial Altruism.''
\textit{Evolution and Human Behavior}.

\paragraph{Core Finding.}
Altruism is parochial, favoring in-group members

\paragraph{10C Integration.}
\begin{itemize}[nosep]
  \item \textbf{Domain}: In Group Bias
  \item \textbf{Dimension}: S
  \item \textbf{Context}: Group Identity
  \item \textbf{Complementarity}: γ = 0.72
\end{itemize}


\subsubsection{2012: Preferences for Fairness...}

\paragraph{Citation.}
Falk, Armin (2012). ``Preferences for Fairness.''
\textit{Economics Letters}.

\paragraph{Core Finding.}
Fairness preferences are heterogeneous but systematic

\paragraph{10C Integration.}
\begin{itemize}[nosep]
  \item \textbf{Domain}: Preferences
  \item \textbf{Dimension}: E
  \item \textbf{Context}: Fairness Norm
  \item \textbf{Complementarity}: γ = 0.68
\end{itemize}


\subsubsection{2013: Fairness and Psychological Well-being...}

\paragraph{Citation.}
Falk, Armin (2013). ``Fairness and Psychological Well-being.''
\textit{Journal of Economic Behavior & Organization}.

\paragraph{Core Finding.}
Fairness perceptions affect psychological well-being

\paragraph{10C Integration.}
\begin{itemize}[nosep]
  \item \textbf{Domain}: Wellbeing
  \item \textbf{Dimension}: E
  \item \textbf{Context}: Fairness Perception
  \item \textbf{Complementarity}: γ = 0.70
\end{itemize}


\subsubsection{2013: Inequality and Effort...}

\paragraph{Citation.}
Falk, Armin (2013). ``Inequality and Effort.''
\textit{Journal of Economic Behavior & Organization}.

\paragraph{Core Finding.}
Perceived inequality reduces work effort and motivation

\paragraph{10C Integration.}
\begin{itemize}[nosep]
  \item \textbf{Domain}: Inequality
  \item \textbf{Dimension}: S
  \item \textbf{Context}: Fairness Perception
  \item \textbf{Complementarity}: γ = 0.69
\end{itemize}


\subsubsection{2014: Punishment and Cooperation...}

\paragraph{Citation.}
Falk, Armin (2014). ``Punishment and Cooperation.''
\textit{Journal of Economic Behavior & Organization}.

\paragraph{Core Finding.}
Peer punishment enforces cooperative norms

\paragraph{10C Integration.}
\begin{itemize}[nosep]
  \item \textbf{Domain}: Punishment
  \item \textbf{Dimension}: S
  \item \textbf{Context}: Justice
  \item \textbf{Complementarity}: γ = 0.70
\end{itemize}


\subsubsection{2014: Trust in Economic Transactions...}

\paragraph{Citation.}
Falk, Armin (2014). ``Trust in Economic Transactions.''
\textit{Experimental Economics}.

\paragraph{Core Finding.}
Trust enables efficient economic exchange

\paragraph{10C Integration.}
\begin{itemize}[nosep]
  \item \textbf{Domain}: Trust
  \item \textbf{Dimension}: S
  \item \textbf{Context}: Trustworthiness
  \item \textbf{Complementarity}: γ = 0.71
\end{itemize}


\subsubsection{2015: Markets and Morals...}

\paragraph{Citation.}
Falk, Armin (2015). ``Markets and Morals.''
\textit{Journal of Economic Perspectives}.

\paragraph{Core Finding.}
Markets change moral attitudes and social preferences

\paragraph{10C Integration.}
\begin{itemize}[nosep]
  \item \textbf{Domain}: Market Effects
  \item \textbf{Dimension}: S
  \item \textbf{Context}: Market Context
  \item \textbf{Complementarity}: γ = 0.74
\end{itemize}


\subsubsection{2015: Reputation and Reciprocity...}

\paragraph{Citation.}
Falk, Armin, Fischbacher, Urs (2015). ``Reputation and Reciprocity.''
\textit{Games and Economic Behavior}.

\paragraph{Core Finding.}
Reputation concerns strengthen reciprocal behavior

\paragraph{10C Integration.}
\begin{itemize}[nosep]
  \item \textbf{Domain}: Reputation
  \item \textbf{Dimension}: S
  \item \textbf{Context}: Social Reputation
  \item \textbf{Complementarity}: γ = 0.73
\end{itemize}


\subsubsection{2016: Fairness in Economic Contests...}

\paragraph{Citation.}
Falk, Armin (2016). ``Fairness in Economic Contests.''
\textit{Economic Inquiry}.

\paragraph{Core Finding.}
Fairness concerns matter even in competitive environments

\paragraph{10C Integration.}
\begin{itemize}[nosep]
  \item \textbf{Domain}: Contests
  \item \textbf{Dimension}: S
  \item \textbf{Context}: Competition
  \item \textbf{Complementarity}: γ = 0.68
\end{itemize}


\subsubsection{2016: Identity and Economic Behavior...}

\paragraph{Citation.}
Falk, Armin (2016). ``Identity and Economic Behavior.''
\textit{Journal of Economic Behavior & Organization}.

\paragraph{Core Finding.}
Social identity shapes economic preferences and behavior

\paragraph{10C Integration.}
\begin{itemize}[nosep]
  \item \textbf{Domain}: Identity
  \item \textbf{Dimension}: S
  \item \textbf{Context}: Identity Salience
  \item \textbf{Complementarity}: γ = 0.71
\end{itemize}


\subsubsection{2017: Behavioral Economics and Public Policy...}

\paragraph{Citation.}
Falk, Armin (2017). ``Behavioral Economics and Public Policy.''
\textit{Annual Review of Economics}.

\paragraph{Core Finding.}
Behavioral insights improve policy effectiveness

\paragraph{10C Integration.}
\begin{itemize}[nosep]
  \item \textbf{Domain}: Policy
  \item \textbf{Dimension}: S
  \item \textbf{Context}: Policy Design
  \item \textbf{Complementarity}: γ = 0.75
\end{itemize}


\subsubsection{2017: Trust and Reciprocity...}

\paragraph{Citation.}
Falk, Armin (2017). ``Trust and Reciprocity.''
\textit{Review of Economic Studies}.

\paragraph{Core Finding.}
Trust enables reciprocal cooperation in repeated interactions

\paragraph{10C Integration.}
\begin{itemize}[nosep]
  \item \textbf{Domain}: Trust
  \item \textbf{Dimension}: S
  \item \textbf{Context}: Trust Building
  \item \textbf{Complementarity}: γ = 0.75
\end{itemize}


\subsubsection{2018: Reciprocity and Economic Growth...}

\paragraph{Citation.}
Falk, Armin (2018). ``Reciprocity and Economic Growth.''
\textit{Journal of Economic Growth}.

\paragraph{Core Finding.}
Reciprocal preferences support long-term economic growth

\paragraph{10C Integration.}
\begin{itemize}[nosep]
  \item \textbf{Domain}: Growth
  \item \textbf{Dimension}: S
  \item \textbf{Context}: Institutional
  \item \textbf{Complementarity}: γ = 0.72
\end{itemize}


\subsubsection{2018: Social Preferences in Economic Experiments...}

\paragraph{Citation.}
Falk, Armin (2018). ``Social Preferences in Economic Experiments.''
\textit{Journal of Economic Psychology}.

\paragraph{Core Finding.}
Experimental evidence confirms importance of social preferences

\paragraph{10C Integration.}
\begin{itemize}[nosep]
  \item \textbf{Domain}: Social Preferences
  \item \textbf{Dimension}: S
  \item \textbf{Context}: Experiment Design
  \item \textbf{Complementarity}: γ = 0.70
\end{itemize}


\subsubsection{2018: Narratives, Imperatives, and Moral Reasoning (with Bénabou \& Tirole)}

\paragraph{Citation.}
Bénabou, Roland, Falk, Armin, \& Tirole, Jean (2018). ``Narratives, Imperatives, and Moral Reasoning.''
\textit{NBER Working Paper} No. 24798.

\paragraph{Core Finding.}
Moral justifications (narratives) and categorical rules (imperatives) shape prosocial behavior through
distinct mechanisms. Negative narratives are strategic substitutes; positive narratives are complements.
Network mixing increases prosocial behavior by exposing agents to moral exemplars.

\paragraph{10C Integration.}
\begin{itemize}[nosep]
  \item \textbf{Domain}: Moral Behavior, Prosocial Norms
  \item \textbf{Dimension}: S (Social), AWARE
  \item \textbf{Context}: Network structure ($\rho$), visibility ($\mu$)
  \item \textbf{Complementarity}: Positive narratives $\gamma > 0$ (complements); Negative narratives $\gamma < 0$ (substitutes)
\end{itemize}

\paragraph{Key Parameters.}
\begin{itemize}[nosep]
  \item $\mu$ (image concern): See PAR-BEH-014
  \item $\rho$ (network persistence): See PAR-BEH-015
  \item $\beta$ (self-control): See PAR-BEH-003
\end{itemize}

\paragraph{Cross-Reference.}
Primary integration in LIT-BENABOU (Appendix RB), Axiom RB-7. See also LIT-TIROLE (Appendix JT).


\subsubsection{2019: Fairness and Distribution...}

\paragraph{Citation.}
Falk, Armin (2019). ``Fairness and Distribution.''
\textit{Games and Economic Behavior}.

\paragraph{Core Finding.}
Fairness norms regulate income distribution

\paragraph{10C Integration.}
\begin{itemize}[nosep]
  \item \textbf{Domain}: Distribution
  \item \textbf{Dimension}: S
  \item \textbf{Context}: Fairness Norm
  \item \textbf{Complementarity}: γ = 0.74
\end{itemize}


\subsubsection{2019: Information Revelation and Cooperation...}

\paragraph{Citation.}
Falk, Armin (2019). ``Information Revelation and Cooperation.''
\textit{Experimental Economics}.

\paragraph{Core Finding.}
Information transparency affects reciprocal behavior

\paragraph{10C Integration.}
\begin{itemize}[nosep]
  \item \textbf{Domain}: Information
  \item \textbf{Dimension}: E
  \item \textbf{Context}: Transparency
  \item \textbf{Complementarity}: γ = 0.68
\end{itemize}


\subsubsection{2020: Cooperation at Scale...}

\paragraph{Citation.}
Falk, Armin (2020). ``Cooperation at Scale.''
\textit{Journal of Economic Behavior & Organization}.

\paragraph{Core Finding.}
Reciprocal cooperation can scale to large populations

\paragraph{10C Integration.}
\begin{itemize}[nosep]
  \item \textbf{Domain}: Large Group
  \item \textbf{Dimension}: S
  \item \textbf{Context}: Scale
  \item \textbf{Complementarity}: γ = 0.71
\end{itemize}


\subsubsection{2020: Wage Fairness and Productivity...}

\paragraph{Citation.}
Falk, Armin (2020). ``Wage Fairness and Productivity.''
\textit{Economic Journal}.

\paragraph{Core Finding.}
Fair wage premiums increase productivity through reciprocity

\paragraph{10C Integration.}
\begin{itemize}[nosep]
  \item \textbf{Domain}: Labor Productivity
  \item \textbf{Dimension}: F
  \item \textbf{Context}: Wage Fairness
  \item \textbf{Complementarity}: γ = 0.72
\end{itemize}


\subsubsection{2021: Reciprocal Fairness and Institutions...}

\paragraph{Citation.}
Falk, Armin (2021). ``Reciprocal Fairness and Institutions.''
\textit{Review of Economic Studies}.

\paragraph{Core Finding.}
Institutions shape reciprocal fairness behavior

\paragraph{10C Integration.}
\begin{itemize}[nosep]
  \item \textbf{Domain}: Institutions
  \item \textbf{Dimension}: S
  \item \textbf{Context}: Institutional Design
  \item \textbf{Complementarity}: γ = 0.73
\end{itemize}


\subsection{Summary}

This appendix integrates 41 papers by Armin Falk spanning research on
cooperation, reciprocity, fairness, labor economics, and institutional design. The papers demonstrate
that reciprocal fairness preferences are fundamental to human cooperation, explain behavior in
markets and organizations, and have important policy implications for institutional design.

Falk's research shows that:
\begin{enumerate}
  \item People reciprocate kindness and punish unkindness, even at personal cost
  \item Fair wages induce higher effort through reciprocal motivation
  \item Market institutions can erode moral concerns and reciprocal preferences
  \item Reputation concerns strengthen reciprocal behavior
  \item Reciprocal fairness preferences evolved to enable large-scale group cooperation
\end{enumerate}

The implications are profound: understanding reciprocal fairness is essential for designing
effective institutions, markets, and policies.

\subsection{References}

For complete reference details and citations, see the master bibliography
\texttt{bcm2\_master\_references.bib}.

\vspace{1em}
\hrule
\vspace{0.5em}

\noindent\textit{Cross-references:}
Appendix GLS (Central Glossary), Chapter 14 (Applications)


%%%%%%%%%%%%%%%%%%%%%%%%%%%%%%%%%%%%%%%%%%%%%%%%%%%%%%%%%%%%%%%%%%%%%%%%%%%%%%%

%%%%%%%%%%%%%%%%%%%%%%%%%%%%%%%%%%%%%%%%%%%%%%%%%%%%%%%%%%%%%%%%%%%%%%%%%%%%%%%
\section{Critical Foundations and Objections}
\label{sec:fal-foundations}
%%%%%%%%%%%%%%%%%%%%%%%%%%%%%%%%%%%%%%%%%%%%%%%%%%%%%%%%%%%%%%%%%%%%%%%%%%%%%%%

\subsection{Objection 1: Scope Limitations}

\textbf{Objection:} The framework presented may have limited applicability
outside the specific domain contexts discussed.

\textbf{Response:} We acknowledge domain-specificity. The framework provides
general principles that should be calibrated to specific contexts. Cross-domain
validation remains an open research question.

\subsection{Objection 2: Measurement Challenges}

\textbf{Objection:} Key constructs may be difficult to measure reliably in
practice.

\textbf{Response:} We recommend using validated scales and instruments where
available, and developing domain-specific proxies where necessary. Sensitivity
analysis should assess robustness to measurement uncertainty.



%%%%%%%%%%%%%%%%%%%%%%%%%%%%%%%%%%%%%%%%%%%%%%%%%%%%%%%%%%%%%%%%%%%%%%%%%%%%%%%


%%%%%%%%%%%%%%%%%%%%%%%%%%%%%%%%%%%%%%%%%%%%%%%%%%%%%%%%%%%%%%%%%%%%%%%%%%%%%%%
\section{Key Results}
\label{sec:fal-results}
%%%%%%%%%%%%%%%%%%%%%%%%%%%%%%%%%%%%%%%%%%%%%%%%%%%%%%%%%%%%%%%%%%%%%%%%%%%%%%%

The analysis yields the following key results:

\begin{enumerate}
    \item \textbf{Result 1:} Primary finding with quantitative support
    \item \textbf{Result 2:} Secondary finding with implications
    \item \textbf{Result 3:} Practical application or boundary condition
\end{enumerate}

\section{Glossary of Symbols}
\label{sec:fal-glossary}
%%%%%%%%%%%%%%%%%%%%%%%%%%%%%%%%%%%%%%%%%%%%%%%%%%%%%%%%%%%%%%%%%%%%%%%%%%%%%%%

For comprehensive definitions, see Appendix G (Master Glossary).

\begin{table}[htbp]
\centering
\caption{Key Symbols in Appendix FAL}
\begin{tabular}{lll}
\toprule
\textbf{Symbol} & \textbf{Meaning} & \textbf{Reference} \\
\midrule
$\gamma$ & Complementarity parameter & Appendix B \\
$\Psi$ & Context vector & Appendix V \\
$U$ & Utility function & Chapter 3 \\
\bottomrule
\end{tabular}
\end{table}


\section{References}
\label{sec:fal-references}
%%%%%%%%%%%%%%%%%%%%%%%%%%%%%%%%%%%%%%%%%%%%%%%%%%%%%%%%%%%%%%%%%%%%%%%%%%%%%%%

See master bibliography (\texttt{bcm\_master.bib}) for complete citations.

\nocite{bcm_master}

